\documentclass[../DoAn.tex]{subfiles}
\begin{document}
\sloppy

% ===== ĐỊNH NGHĨA TỪ KHÓA THUẬT TOÁN CHO CHƯƠNG 5 =====
\SetKwIF{If}{ElseIf}{Else}{IF}{THEN}{ELSE IF}{ELSE}{END IF}
\SetKwFor{For}{FOR}{DO}{END FOR}
\SetKwFor{While}{WHILE}{DO}{END WHILE}
\SetKwSwitch{Switch}{Case}{Other}{CASE}{OF}{}{}{END CASE}{}
\SetKw{Begin}{BEGIN}
\SetKw{End}{END}
\SetKw{KwReturn}{RETURN}
\SetKw{KwBreak}{BREAK}
\SetKw{KwCall}{CALL}
\SetKw{KwException}{EXCEPTION}
\SetAlgoNoLine  % Bỏ đường kẻ ngang trong tất cả thuật toán
\SetAlgoSkip{}          % bỏ skip trên/dưới algorithm
\setlength{\algomargin}{1em}  % giảm indent margin
% =======================================================

Bốn chương trước đã trình bày bài toán, yêu cầu, công nghệ và thiết kế của hệ
thống theo hướng hệ thống hóa. Trên cơ sở đó, chương này tập trung phân tích
ba đóng góp kỹ thuật chính của đồ án, không chỉ ở mức mô tả chức năng mà còn ở
cách đồ án đã thiết kế, thử nghiệm và điều chỉnh giải pháp để đáp ứng các yêu
cầu thực tế của bài toán quản lý ký túc xá.

Mỗi đóng góp vẫn trả lời bốn câu hỏi cốt lõi: \textit{(i)} tính năng giải
quyết vấn đề gì; \textit{(ii)} đối tượng nào được hưởng lợi; \textit{(iii)}
điểm khác biệt so với giải pháp hiện có là gì; \textit{(iv)} giá trị mang lại
cho hệ thống và người dùng ra sao. Chương được tổ chức thành bốn mục: giải
pháp trực quan hóa không gian phòng ký túc xá~(5.1), giải pháp phân bổ phòng
tự động dựa trên quota~(5.2), giải pháp kiến trúc phân bổ linh hoạt theo năm
học~(5.3) và tổng kết chương~(5.4).

% =====================================================================
\section{Giải pháp trực quan hóa không gian phòng ký túc xá}
\label{sec:contrib_visualization}
% =====================================================================

\subsection{Dẫn dắt và bài toán}
\label{subsec:viz_problem}

\subsubsection{Vấn đề cần giải quyết}

Phần lớn sinh viên nhập học tại HUST đến từ các tỉnh thành ngoài Hà Nội và
chưa từng đặt chân đến ký túc xá trước khi đăng ký. Hệ thống hiện tại chỉ
cung cấp ảnh tĩnh và mô tả văn bản --- không đủ để hình thành nhận thức không
gian thực tế, vì mắt người cần tự do chuyển hướng nhìn mới cảm nhận được
chiều sâu và tỷ lệ. Hệ quả là một tỷ lệ đáng kể sinh viên phát hiện phòng
không khớp kỳ vọng sau khi nhận, dẫn đến yêu cầu đổi phòng và tăng tải công
việc cho cán bộ quản lý.

Bài toán đặt ra không chỉ là bổ sung một trang xem ảnh phòng, mà là thiết kế
module trực quan hóa đủ chân thực để sinh viên có căn cứ quan sát trước khi
đăng ký từ xa, đồng thời vẫn đảm bảo hiệu suất tải hợp lý trên web và khả dụng
trên cả thiết bị di động.

\subsubsection{Đối tượng hưởng lợi}

Sinh viên đăng ký có thể đi vào phòng, nhìn quanh 360°, di
chuyển tự do trong không gian 3D trước khi nhấn nút đăng ký --- quyết định
dựa trên căn cứ trực quan thay vì suy đoán hoặc thông tin truyền miệng. Cán bộ
quản lý giảm số yêu cầu đổi phòng sau đăng ký và có công cụ tự cập nhật nội
dung trực quan mà không cần lập trình viên. Hệ thống E-Dormitory hưởng lợi qua
chất lượng dữ liệu: lựa chọn phòng của sinh viên phản ánh đúng sở thích hơn,
giảm nhiễu cho thuật toán phân bổ.

\subsubsection{Điểm khác biệt so với giải pháp hiện có}

Hầu hết ký túc xá đại học tại Việt Nam dừng lại ở ảnh tĩnh hoặc video tour
tuyến tính --- cả hai đều chỉ cho phép xem thụ động. Điểm khác biệt then chốt
của giải pháp đề xuất là chuyển dịch từ xem thụ động sang trải nghiệm
chủ động: người dùng tự quyết định nhìn vào đâu và đứng ở đâu trong phòng.
Bảng~\ref{tab:viz_comparison} tổng hợp so sánh ba phương án.

\begin{table}[H]
\centering
\caption{So sánh các phương án trực quan hóa không gian phòng}
\label{tab:viz_comparison}
\small
\begin{tabular}{|p{2.8cm}|p{2.6cm}|p{2.6cm}|p{5.2cm}|}
\hline
\textbf{Tiêu chí} & \textbf{Ảnh tĩnh} & \textbf{Video tour} & \textbf{Giải pháp đề xuất (360° + 3D)} \\
\hline
Nhận thức không gian & Rất hạn chế & Trung bình & Cao --- cảm nhận chiều sâu và tỷ lệ thực \\
\hline
Tính tương tác & Không có & Thụ động & Chủ động --- người dùng tự chọn góc nhìn \\
\hline
Di chuyển tự do & Không & Không & Có (chế độ Free Explore trên desktop) \\
\hline
Tương thích thiết bị & Mọi thiết bị & Mọi thiết bị & 360° mọi thiết bị; 3D trên desktop \\
\hline
Cập nhật nội dung & Cần kỹ thuật & Cần kỹ thuật & Cán bộ tự cập nhật qua giao diện \\
\hline
Thời gian sẵn sàng & Ngay lập tức & Phụ thuộc băng thông & < 0,8 giây đến tương tác đầu tiên \\
\hline
\end{tabular}
\end{table}

\subsection{Giải pháp}
\label{subsec:viz_solution}

\subsubsection{Hành trình thử nghiệm --- từ ảnh tĩnh đến mô hình 3D}

Đây là đóng góp đòi hỏi nhiều vòng thử nghiệm và điều chỉnh nhất trong quá
trình thực hiện đồ án. Giải pháp cuối cùng không xuất hiện ngay từ đầu mà được
hình thành qua nhiều bước: mỗi phương án ban đầu giải quyết được một phần vấn
đề, đồng thời làm lộ ra một hạn chế mới cần xử lý. Hành trình đó trải qua bốn
giai đoạn, mỗi giai đoạn giải quyết một hạn chế của giai đoạn trước, được minh
họa trong Hình~\ref{fig:viz_journey}.

\begin{figure}[H]
\centering
\includegraphics[width=0.98\textwidth]{Hinhve/quy_trinh_xay_dung_3d.png}
\caption{Hành trình bốn giai đoạn xây dựng giải pháp trực quan hóa}
\label{fig:viz_journey}
\end{figure}

Trong quá trình xây dựng chức năng trực quan hóa phòng ký túc xá, đồ án bắt đầu
từ giải pháp đơn giản nhất là sử dụng nhiều ảnh chụp tĩnh. Mỗi phòng được chụp từ
nhiều góc khác nhau, kèm theo thông tin về kích thước và vị trí của từng khu
vực nhằm giúp sinh viên có thể hình dung bố cục không gian trước khi đăng ký.

Sau khi thử nghiệm với người dùng, cách tiếp cận này bộc lộ nhiều hạn chế.
Mặc dù số lượng ảnh khá lớn, người xem vẫn khó ước lượng được chiều sâu và tỷ
lệ thực của căn phòng. Nguyên nhân là ảnh tĩnh chỉ cung cấp các góc nhìn rời
rạc, trong khi nhận thức không gian của con người phụ thuộc nhiều vào khả năng
thay đổi góc nhìn liên tục. Điều này cho thấy vấn đề không nằm ở số lượng hình
ảnh mà nằm ở việc người dùng chưa thể chủ động quan sát môi trường.

Để khắc phục hạn chế trên, đồ án chuyển sang sử dụng ảnh toàn cảnh
equirectangular 360° kết hợp với thư viện Photo Sphere
Viewer~\cite{psv2023}. Thay vì chỉ xem các ảnh riêng lẻ, người dùng có thể kéo
chuột để quan sát toàn bộ không gian từ một vị trí chụp cố định. Trải nghiệm
này giúp cảm nhận rõ hơn về chiều sâu, khoảng cách và bố cục của căn phòng.

Tuy nhiên, sau khi hoàn thành phiên bản đầu tiên, quá trình kiểm thử cho thấy
một hạn chế khác: người dùng chỉ có thể quan sát từ đúng vị trí đặt máy ảnh. Trong thực
tế, một phòng ký túc xá có nhiều khu vực quan trọng như giường ngủ, bàn học,
tủ đồ hay khu vệ sinh, trong khi mỗi vị trí chỉ thể hiện tốt một phần của
không gian. Vì vậy, đồ án bổ sung hệ thống hotspot cho phép người dùng chuyển
đổi giữa nhiều điểm quan sát khác nhau ngay trong ảnh 360°.

Việc triển khai hotspot phát sinh thêm một bài toán kỹ thuật. Photo Sphere
Viewer sử dụng hệ tọa độ cầu, trong khi vị trí các điểm trong mô hình 3D được
biểu diễn bằng hệ tọa độ Descartes của Three.js. Do đó, mỗi hotspot phải được
chuyển đổi từ tọa độ $(x,y,z)$ sang cặp góc $(\theta,\phi)$ theo công thức

\begin{equation}
\theta = \arctan!\left(\frac{x}{-z}\right), \quad
\phi = \arctan!\left(\frac{y}{\sqrt{x^2+z^2}}\right).
\label{eq:cart_to_sphere}
\end{equation}

Trong quá trình triển khai thực tế, việc xác định các giá trị
$(\theta,\phi)$ thủ công cho thấy không phù hợp với cán bộ quản lý ký túc xá vì họ không
có kiến thức lập trình. Để giải quyết vấn đề này, đồ án xây dựng một công cụ
tích hợp ngay trong giao diện quản trị. Người quản trị chỉ cần di chuyển chuột
đến vị trí mong muốn trên ảnh 360°, nhấn nút lấy tọa độ và hệ thống sẽ tự động
tính toán, sau đó lưu các giá trị tương ứng vào cơ sở dữ liệu. Nhờ vậy, việc
thêm hoặc cập nhật hotspot cho các phòng mới không còn yêu cầu chỉnh sửa mã
nguồn.

Mặc dù chế độ 360° đã cải thiện đáng kể khả năng quan sát, người dùng vẫn chỉ
có thể di chuyển giữa các vị trí được định nghĩa sẵn. Vì vậy, đồ án tiếp tục
xây dựng một chế độ cho phép người dùng tự do di chuyển trong phòng giống như
đang trực tiếp đứng bên trong không gian đó. Do không có sẵn mô hình số của
các phòng ký túc xá HUST, sinh viên thực hiện đồ án phải tự học Blender~\cite{blender2024}, đo đạc
kích thước thực tế và dựng lại toàn bộ mô hình từ đầu. Công việc bao gồm dựng
tường, sàn, trần, cửa, cửa sổ, cùng toàn bộ nội thất như giường tầng, bàn học,
tủ cá nhân và các vật dụng khác.

Sau khi hoàn thành mô hình đầu tiên, mô hình được xuất sang định dạng GLB để tích hợp
lên hệ thống. Tuy nhiên, kích thước tệp vượt quá 200~MB, lớn hơn rất nhiều so
với mức có thể sử dụng trong một ứng dụng web. Thời gian tải kéo dài, lượng bộ
nhớ tiêu thụ lớn và tốc độ hiển thị không đáp ứng được yêu cầu thực tế. Thay
vì tiếp tục tối ưu một cách thử nghiệm, đồ án tiến hành phân tích nguyên nhân gây
ra kích thước lớn của mô hình.

Kết quả phân tích cho thấy dung lượng chủ yếu đến từ ba thành phần: số lượng
đa giác quá cao, texture PBR có độ phân giải lớn và dữ liệu hình học chưa được
nén.

Trước hết, đồ án sử dụng Decimate Modifier của Blender để giảm số lượng đa giác
trên các bề mặt phẳng như tường, sàn và trần, đồng thời giữ nguyên những đối
tượng cần nhiều chi tiết như giường hoặc bàn học. Sau bước này, số lượng đa
giác giảm từ hơn 400.000 xuống còn khoảng 120.000.

Tiếp theo, đồ án thực hiện texture baking bằng cách gộp nhiều texture map như
albedo, roughness, metallic và normal thành một texture duy nhất cho từng nhóm
đối tượng. Cách làm này vừa giảm số lượng texture cần tải, vừa giảm số lượng
draw call trong quá trình kết xuất, đưa kích thước mô hình xuống còn khoảng
30--40~MB.

Cuối cùng, đồ án áp dụng thuật toán nén hình học Draco của
Google~\cite{google2023draco}. Draco nén dữ liệu hình học của mô hình nhưng
vẫn giữ nguyên hình dạng sau khi giải nén ở phía trình duyệt. Sau bước tối ưu
này, kích thước tệp GLB giảm xuống dưới 8~MB, tức nhỏ hơn khoảng 25 lần so với
phiên bản ban đầu. Mức dung lượng này cho phép mô hình được tải nhanh trên môi
trường web và đáp ứng yêu cầu triển khai trong hệ thống quản lý ký túc xá.


Hành trình tối ưu hóa được tổng kết trong Bảng~\ref{tab:glb_optimization}.

\begin{table}[H]
\centering
\caption{Tiến trình tối ưu hóa kích thước file GLB}
\label{tab:glb_optimization}
\small
\begin{tabular}{|p{4.2cm}|p{2.2cm}|p{2.5cm}|p{4.3cm}|}
\hline
\textbf{Giai đoạn} & \textbf{Kích thước} & \textbf{Số đa giác} & \textbf{Kỹ thuật áp dụng} \\
\hline
Mô hình gốc (chưa tối ưu) & $\sim$200~MB & $>$400.000 & --- \\
\hline
Sau Decimate & $\sim$35~MB & 120.000 & Hợp nhất đa giác bề mặt phẳng \\
\hline
Sau Texture Baking & $\sim$30~MB & 120.000 & Gộp vật liệu thành texture map \\
\hline
Sau nén Draco & < 8~MB & 22.000--25.000 & Nén hình học + geometry index \\
\hline
\end{tabular}
\end{table}

Ngoài vấn đề kích thước file, còn một thách thức chỉ lộ ra khi cầm điện thoại
lên và thực sự mở trang: trình duyệt không crash ngay --- nó load được nhưng
sau vài giây render thì đóng băng hoàn toàn. Tra log Chrome DevTools mobile
mới phát hiện lỗi thiếu WebGL extension \texttt{OES\allowbreak\_element\allowbreak\_index\allowbreak\_uint}
mà Three.js cần cho index buffer 32-bit khi render mô hình lớn. Thay vì cố
vượt qua giới hạn phần cứng, quyết định thiết kế được đưa ra: ẩn hoàn toàn
chế độ Free Explore trên thiết bị di động và chỉ cung cấp Quick Tour (360°).
Đây là đánh đổi có chủ ý --- tốt hơn là cung cấp trải nghiệm hạn chế nhưng
ổn định, hơn là cố render mô hình và gây crash trình duyệt.

\subsubsection{Kiến trúc hai chế độ song song}

Toàn bộ hành trình bốn giai đoạn hội tụ vào kiến trúc cuối cùng: hai chế độ
song song chia sẻ cùng một tầng dữ liệu, được trình bày trong
Hình~\ref{fig:visualization_architecture}.

\begin{figure}[H]
\centering
\includegraphics[width=0.8\textwidth]{Hinhve/kien_truc_truc_quan_hoa_phong.png}
\caption{Kiến trúc giải pháp trực quan hóa không gian phòng ký túc xá}
\label{fig:visualization_architecture}
\end{figure}

Module tổ chức thành hai nhánh chạy song song. Nhánh \textit{Quick Tour} xử
lý ảnh equirectangular 360° qua Photo Sphere Viewer, khởi tạo ngay lập tức
khi trang tải và khả dụng trên mọi thiết bị. Nhánh \textit{Free Explore} tải
mô hình GLB qua Three.js và chỉ được kích hoạt trên desktop khi mạng đủ nhanh.
Cả hai nhánh chia sẻ tầng siêu dữ liệu phòng và hệ thống hotspot, đảm bảo
tính nhất quán dữ liệu giữa hai chế độ.

Chiến lược tải song song là điểm quan trọng về hiệu suất: thay vì chờ GLB
tải xong mới hiển thị gì đó, Quick Tour được khởi tạo ngay lập tức trong khi
GLB tải ngầm --- giảm thời gian đến tương tác đầu tiên từ hàng giây xuống
dưới 0,8 giây. Bảng~\ref{tab:viz_modes} tóm tắt đặc điểm từng chế độ.

\begin{table}[H]
\centering
\caption{So sánh hai chế độ trực quan hóa}
\label{tab:viz_modes}
\small
\begin{tabular}{|p{2.8cm}|p{5.4cm}|p{5.0cm}|}
\hline
\textbf{Tiêu chí} & \textbf{Quick Tour (360°)} & \textbf{Free Explore (3D)} \\
\hline
Công nghệ & Photo Sphere Viewer + Three.js & Three.js WebGL + GLB \\
\hline
Điều hướng & Hotspot tọa độ cầu xác định trước & Di chuyển tự do WASD + chuột \\
\hline
Tải trọng & Thấp (JPEG tối ưu) & < 8~MB (GLB nén Draco) \\
\hline
Thiết bị hỗ trợ & Desktop + mobile + tablet & Desktop (WebGL đầy đủ) \\
\hline
Thời gian sẵn sàng & Ngay lập tức & Sau 1,2--1,8 giây (tải ngầm) \\
\hline
\end{tabular}
\end{table}

\subsubsection{Thuật toán tối ưu hóa mô hình 3D}

\paragraph{Mục tiêu.}
Giảm kích thước file GLB từ mô hình gốc (hàng trăm MB) xuống mức chấp nhận
được cho ứng dụng web (dưới 10~MB) mà không làm mất chất lượng hình ảnh cảm
nhận được bởi người dùng, thông qua ba kỹ thuật áp dụng tuần tự.

\begin{algorithm}[H]
\small
\caption{\textit{GLBModelOptimizationPipeline}}
\KwData{$RawModel$: mô hình Blender gốc; $TargetSize$: kích thước mục tiêu (MB);
        $QualityThreshold$: ngưỡng chất lượng chấp nhận được}
\KwResult{$OptimizedGLB$: file GLB đã tối ưu; $Metrics$: báo cáo kích thước và chất lượng}
\tcp{Bước 1: Decimate}
\For{$obj \in RawModel.objects$}{
  \If{$obj.type \in \{\text{FLOOR, WALL, CEILING}\}$}{
    $obj.decimateRatio \leftarrow 0.05$\; \tcp*[r]{Bề mặt phẳng --- giảm mạnh}
  }
  \Else{
    $obj.decimateRatio \leftarrow 0.4$\; \tcp*[r]{Đồ nội thất --- giữ chi tiết}
  }
  $\text{ApplyDecimate}(obj,\ obj.decimateRatio)$\;
}
\tcp{Bước 2: Texture Baking}
\For{$group \in \text{GroupByMaterial}(RawModel.objects)$}{
  $texMap \leftarrow \text{BakeToTexture}(group,\ \text{resolution:}\ 1024\times1024)$\;
  $\text{ReplaceMaterials}(group,\ texMap)$\;
}
\tcp{Bước 3: Nén Draco --- tìm mức nén tối ưu}
$rawGLB \leftarrow \text{ExportGLB}(RawModel)$\;
$level \leftarrow 7$\;
\While{$rawGLB.size > TargetSize$ \textbf{AND} $level \leq 10$}{
  $candidate \leftarrow \text{ApplyDraco}(rawGLB,\ \text{compressionLevel:}\ level)$\;
  \If{$\text{VisualQuality}(candidate) \geq QualityThreshold$}{
    $OptimizedGLB \leftarrow candidate$\;
    $level \leftarrow level + 1$\;
  }
  \Else{
    \KwBreak\; \tcp*[r]{Vượt ngưỡng chất lượng --- dừng tăng nén}
  }
}
$Metrics \leftarrow \{size{:}\ OptimizedGLB.size,\ polygons{:}\ \text{CountPolygons}(OptimizedGLB)\}$\;
\KwReturn $OptimizedGLB$, $Metrics$\;
\end{algorithm}

\subsubsection{Cơ chế khởi tạo và điều phối chế độ trực quan hóa}

Để giảm thời gian chờ của người dùng trong khi vẫn hỗ trợ chế độ quan sát 3D,
hệ thống được thiết kế theo cơ chế khởi tạo nhiều mức thay vì tải toàn bộ dữ
liệu ngay từ đầu. Khi người dùng mở trang xem phòng, hệ thống trước tiên đọc
metadata của phòng từ cơ sở dữ liệu để xác định các tài nguyên sẵn có, bao gồm
ảnh 360°, danh sách hotspot và mô hình 3D nếu phòng hỗ trợ chế độ Free
Explore.

Ngay sau khi hoàn tất bước khởi tạo, chế độ Quick Tour được ưu tiên tải trước.
Do ảnh equirectangular có dung lượng nhỏ và không yêu cầu xử lý hình học phức
tạp, người dùng có thể bắt đầu quan sát không gian gần như ngay lập tức mà
không phải chờ mô hình 3D hoàn tất.

Trong khi người dùng đang sử dụng Quick Tour, hệ thống tiếp tục đánh giá đặc
điểm thiết bị và điều kiện mạng. Nếu phát hiện thiết bị có đủ năng lực xử lý và
kết nối mạng đáp ứng yêu cầu, mô hình GLB sẽ được tải nền (background loading)
mà không làm gián đoạn quá trình tương tác hiện tại. Khi quá trình tải hoàn
thành, hệ thống tự động kích hoạt chế độ Free Explore và cho phép người dùng
chuyển đổi giữa hai chế độ quan sát bất kỳ lúc nào.

Đối với các thiết bị di động hoặc những trường hợp băng thông thấp, hệ thống
không tải mô hình 3D nhằm tránh hiện tượng thời gian chờ kéo dài hoặc tiêu tốn
quá nhiều bộ nhớ đồ họa. Trong trường hợp này, Quick Tour trở thành chế độ trực
quan hóa chính, bảo đảm người dùng vẫn có đầy đủ khả năng quan sát căn phòng
mà không ảnh hưởng đến trải nghiệm sử dụng.

Sau khi quá trình khởi tạo hoàn tất, toàn bộ thao tác của người dùng đều được
xử lý theo cơ chế hướng sự kiện. Các sự kiện như xoay góc nhìn, lựa chọn
hotspot hoặc chuyển đổi giữa Quick Tour và Free Explore chỉ cập nhật trạng thái
hiển thị hiện tại mà không cần tải lại toàn bộ dữ liệu. Nhờ đó, mọi thao tác
đều được thực hiện gần như tức thời, đồng thời hạn chế việc sử dụng tài nguyên
không cần thiết.

\subsection{Giá trị mang lại}
\label{subsec:viz_value}

Giá trị chính của đóng góp này là chuyển quá trình tìm hiểu phòng từ chỗ phụ
thuộc nhiều vào mô tả ngắn, ảnh rời rạc hoặc truyền miệng sang một trải nghiệm
có căn cứ trực quan hơn. Sinh viên có thể quan sát bố cục, khoảng cách giữa
giường, bàn học, tủ đồ và khu vệ sinh trước khi đăng ký, nhờ đó kỳ vọng về
phòng ở được hình thành dựa trên hình ảnh không gian cụ thể hơn.

Đối với cán bộ quản lý, module trực quan hóa trở thành công cụ hỗ trợ giới
thiệu phòng và ký túc xá rõ ràng hơn trong các đợt tư vấn tuyển sinh hoặc tiếp
nhận sinh viên mới. Thay vì giải thích hoàn toàn bằng lời hoặc gửi một vài ảnh
tĩnh, cán bộ có thể dẫn sinh viên xem từng khu vực trong phòng, chỉ ra sự khác
biệt giữa các loại phòng và cập nhật nội dung tour khi có thay đổi thực tế.

Về phía hệ thống, việc sinh viên có thêm căn cứ quan sát trước khi đăng ký giúp
giảm khả năng chọn phòng theo thông tin không đầy đủ. Điều này không loại bỏ
hoàn toàn nhu cầu đổi phòng, nhưng góp phần giảm các yêu cầu phát sinh do kỳ
vọng ban đầu không khớp với không gian thực tế, từ đó giảm tải xử lý cho cán bộ
và làm dữ liệu đăng ký phản ánh nhu cầu thực tế hơn.




% =====================================================================
\section{Giải pháp phân bổ phòng tự động dựa trên quota}
\label{sec:contrib_allocation}
% =====================================================================

\subsection{Dẫn dắt và bài toán}
\label{subsec:alloc_problem}

\subsubsection{Vấn đề cần giải quyết}

Quá trình trao đổi với cán bộ quản lý cho thấy một đợt phân bổ thủ công
thường mất từ ba đến năm ngày làm việc --- gộp dữ liệu từ nhiều file Excel, lọc
theo giới tính, kiểm tra quota từng nhóm năm và in danh sách để trình ký. Trong
quy trình này, sai sót có thể xuất hiện ở nhiều bước: một phòng được gán cho cả
nam lẫn nữ, nhóm năm nhất vượt quota vì tính sót một cột, hoặc thứ tự danh sách
chờ không còn nhất quán sau nhiều lần chỉnh sửa. Sinh viên cũng khó kiểm chứng
căn cứ vì sao mình được phân phòng hoặc được đưa vào danh sách chờ.

Bài toán phân bổ phòng ký túc xá mang tính ràng buộc đa chiều. Vì vậy, đồ án
không chỉ xây dựng chức năng gán phòng tự động, mà còn tách rõ các ràng buộc
bắt buộc, cơ chế xếp hạng và nhật ký quyết định để mỗi kết quả phân bổ có thể
được kiểm tra lại. Mỗi quyết định gán phòng phải đồng thời thỏa mãn bốn nhóm
ràng buộc --- sức chứa phòng, đồng nhất giới tính, chỉ tiêu theo nhóm năm học
(quota), và thứ tự ưu tiên theo chính sách nhà trường. Đặc biệt, số lượng đơn
đăng ký thường vượt số chỗ ở khả dụng, đòi hỏi một cơ chế xếp hạng minh bạch
và có thể giải trình.

\subsubsection{Đối tượng hưởng lợi}

\textit{Sinh viên đăng ký:} có thể tra cứu lý do cụ thể tại sao mình được
phân phòng hay đưa vào danh sách chờ, thay vì nhận kết quả mà không hiểu cơ
sở. \textit{Cán bộ quản lý:} quy trình phân bổ từ nhiều ngày rút xuống dưới
2 giây, giảm đáng kể thao tác Excel và tính điểm thủ công. \textit{Ban
lãnh đạo ký túc xá:} có nhật ký kiểm toán đầy đủ cho mỗi đợt phân bổ, phục
vụ báo cáo và giải trình khi có khiếu nại.

\subsubsection{Điểm khác biệt so với giải pháp hiện có}

\begin{table}[H]
\centering
\caption{So sánh giải pháp phân bổ phòng}
\label{tab:alloc_comparison}
\small
\begin{tabular}{|p{3.5cm}|p{4.5cm}|p{4.8cm}|}
\hline
\textbf{Tiêu chí} & \textbf{Giải pháp thủ công hiện có} & \textbf{Giải pháp đề xuất} \\
\hline
Thời gian xử lý & Nhiều ngày công tác & Dưới 2 giây \\
\hline
Kiểm tra ràng buộc & Dễ bỏ sót khi làm thủ công & Tự động, tường minh, thứ tự cố định \\
\hline
Tính minh bạch & Sinh viên không biết lý do & Mỗi quyết định có AuditLog riêng \\
\hline
Quản lý danh sách chờ & Thủ công, dễ nhầm thứ tự & Tự động, có thứ tự theo điểm \\
\hline
Khả năng kiểm toán & Không có & Ghi đầy đủ căn cứ từng quyết định \\
\hline
\end{tabular}
\end{table}

\subsection{Giải pháp}
\label{subsec:alloc_model}

\subsubsection{Mô hình hai lớp quyết định}

Ban đầu, toàn bộ logic phân bổ được viết trong một hàm duy nhất: kiểm tra
ràng buộc, tính điểm, gán phòng --- tất cả gộp lại. Khi kiểm thử đồng thời
nhiều đơn, cách tổ chức này bộc lộ rủi ro: hai tiến trình có thể cùng đọc số
chỗ trống và cùng gán phòng, làm kết quả cuối bị sai. Từ đó, đồ án tách
``ràng buộc cứng'' ra kiểm tra trước toàn bộ, đồng thời sử dụng
\texttt{capacitySnapshot} như một bản chụp trạng thái phòng tại thời điểm bắt
đầu phân bổ để hạn chế race condition.

Mô hình phân bổ hoạt động trên hai lớp quyết định độc lập. \textit{Lớp ngoài}
kiểm tra các ràng buộc cứng theo thứ tự cố định: (1)~quota nhóm năm còn chỗ
không; (2)~phòng còn giường trống không; (3)~giới tính phòng có khớp không.
Chỉ khi thỏa mãn cả ba ràng buộc này, ứng viên mới tiếp tục sang lớp trong.
\textit{Lớp trong} sử dụng điểm ưu tiên để xếp hạng các ứng viên cạnh tranh
trong cùng nhóm:
\begin{equation}
Score_i = W_d \cdot D_i + W_f \cdot F_i + W_y \cdot Y_i + W_p \cdot P_i - W_v \cdot V_i
\label{eq:priority_score}
\end{equation}
trong đó $D_i$, $F_i$, $Y_i$, $P_i$, $V_i$ lần lượt là điểm khoảng cách từ
nhà, tài chính khó khăn, nhóm năm học, học bổng và điểm trừ vi phạm nội quy.
Các trọng số $W_\ast$ được lưu trong cấu hình, có thể điều chỉnh mỗi năm mà
không cần sửa mã nguồn. Điểm ưu tiên là \textit{cơ chế xếp hạng}, không phải
cơ chế quyết định: sinh viên phải vượt qua lớp ngoài trước.

Số chỗ ở dành cho từng nhóm năm được tính theo:
\begin{equation}
Quota_g = Capacity \times Ratio_g
\label{eq:quota}
\end{equation}
với $Ratio_g$ là tỷ lệ được cấu hình trong \texttt{QuotaConfig}. Sau mỗi đợt
phân bổ, tỷ lệ lấp đầy được cập nhật trên bảng điều khiển:
\begin{equation}
OccupancyRate = \frac{OccupiedBeds}{TotalBeds} \times 100\%
\label{eq:occupancy}
\end{equation}
Dữ liệu thực tế ghi nhận 82,1\% trên 265 phòng và 1.450 sinh viên.

\subsubsection{Thuật toán phân bổ phòng}

\begin{minipage}{\linewidth}
\paragraph{Mục tiêu.}
Phân bổ toàn bộ danh sách đơn đăng ký vào các phòng khả dụng theo đúng thứ
tự ưu tiên, đảm bảo ba ràng buộc cứng không bao giờ bị vi phạm và mọi quyết
định đều có căn cứ ghi lại trong AuditLog. \\

\begin{algorithm}[H]
\small
\caption{\textit{QuotaBasedRoomAllocation}}
\KwData{$StudentList$: danh sách đơn; $RoomList$: danh sách phòng; $QuotaConfig$: cấu hình quota}
\KwResult{$AssignmentList$: danh sách phân phòng; $WaitingList$: danh sách chờ}
\tcp{Bước 1: Tính điểm và sắp xếp}
\For{$s \in StudentList$}{
  $s.score \leftarrow \text{CalculatePriorityScore}(s)$\;
}
$StudentList \leftarrow \text{SORT}(StudentList,\ key{:}\ score,\ \text{DESC})$\;
\tcp{Bước 2: Phân bổ tuần tự --- kiểm tra 3 ràng buộc cứng}
\For{$s \in StudentList$}{
  \If{$QuotaConfig.slots[s.yearGroup] \leq 0$}{
    $WaitingList.\text{append}(s,\ reason{:}\ \text{QUOTA\_FULL})$\;
    \KwBreak\;
  }
  $room \leftarrow \text{FindRoom}(RoomList,\ gender{:}\ s.gender,\ group{:}\ s.yearGroup)$\;
  \If{$room \neq \varnothing$}{
    $\text{Assign}(s,\ room)$\;
    $QuotaConfig.slots[s.yearGroup] \mathrel{-}= 1$\;
    $\text{WriteAuditLog}(s,\ room,\ \text{ALLOCATED})$\;
  }
  \Else{
    $WaitingList.\text{append}(s,\ reason{:}\ \text{NO\_ROOM})$\;
    $\text{WriteAuditLog}(s,\ \text{WAITLIST})$\;
  }
}
\KwReturn $AssignmentList$, $WaitingList$\;
\end{algorithm}
\end{minipage}

\medskip
\paragraph{Phân tích độ phức tạp.}
Bước tính điểm là $O(n)$; sắp xếp là $O(n \log n)$; phân bổ là $O(n \cdot m)$
trong trường hợp xấu nhất, giảm về $O(n \cdot m/k)$ khi danh sách phòng được
lập chỉ mục theo nhóm năm và giới tính. Với 1.500 đơn và 265 phòng, thuật toán
hoàn thành dưới 2 giây. Tính đúng đắn được đảm bảo bởi $capacitySnapshot$ bất
biến: bản chụp sức chứa ghi lại trước khi phân bổ bắt đầu, tránh race condition
khi cán bộ thao tác phòng đồng thời.

\subsection{Kết quả đạt được}
\label{subsec:alloc_result}

\begin{table}[H]
\centering
\caption{Kết quả kiểm thử xác minh tính đúng đắn thuật toán phân bổ}
\label{tab:alloc_verification}
\small
\begin{tabular}{|p{3.0cm}|p{3.6cm}|p{3.8cm}|p{1.8cm}|}
\hline
\textbf{Kịch bản} & \textbf{Điều kiện đầu vào} & \textbf{Kết quả mong đợi} & \textbf{Kết quả} \\
\hline
Đủ sức chứa & 10 đơn, 15 giường trống & Tất cả 10 đơn ALLOCATED & Đạt \\
\hline
Vượt sức chứa & 20 đơn, 10 giường & 10 điểm cao nhất ALLOCATED, 10 WAITLIST & Đạt \\
\hline
Xếp hạng đúng & 5 đơn điểm khác nhau, 3 giường & 3 đơn điểm cao nhất theo thứ tự giảm dần & Đạt \\
\hline
Hết quota nhóm & 10 đơn năm 1, quota = 5 & 5 ALLOCATED, 5 QUOTA\_FULL & Đạt \\
\hline
Không có giường & Bất kỳ số đơn, 0 giường & Tất cả NO\_ROOM & Đạt \\
\hline
\end{tabular}
\end{table}

So với quy trình thủ công, hệ thống rút ngắn thời gian phân bổ từ vài ngày
xuống dưới 2 giây, giảm thao tác Excel, loại bỏ sai số tính điểm tay và đảm
bảo mọi sinh viên đều có thể tra cứu căn cứ quyết định của mình trong
\texttt{AllocationAuditLog}. Ba ràng buộc cứng được kiểm tra tuần tự và tường
minh trong mã nguồn, không thể bị bỏ qua dù có lỗi ở bước khác.

\subsection{Giá trị mang lại}
\label{subsec:alloc_value}

Giá trị quan trọng nhất của đóng góp này không chỉ nằm ở tốc độ xử lý, mà ở
việc chuyển quy trình phân bổ từ thao tác Excel và điều chỉnh thủ công sang
một quy trình tự động có thứ tự xử lý rõ ràng. Các bước tính điểm, xếp hạng,
kiểm tra quota, tìm phòng và cập nhật kết quả được thực hiện theo cùng một
logic cho mọi đợt phân bổ, nhờ đó kết quả nhất quán hơn giữa các lần chạy.

Cơ chế ghi \texttt{AllocationAuditLog} giúp mỗi quyết định phân bổ có căn cứ
để truy vết: sinh viên nào được xếp hạng ở vị trí nào, được gán phòng hay đưa
vào danh sách chờ vì lý do gì, và đợt phân bổ được thực thi bởi ai. Điều này
tăng tính khách quan và minh bạch của quy trình, đồng thời hỗ trợ cán bộ giải
trình khi có yêu cầu đối chiếu hoặc khiếu nại.

Đối với vận hành ký túc xá, việc tự động hóa không loại bỏ vai trò của cán bộ
quản lý, mà giảm các thao tác lặp lại và dễ sai sót để cán bộ tập trung vào
những trường hợp cần xem xét nghiệp vụ. Kết quả phân bổ vì vậy có căn cứ rõ
ràng hơn, nhất quán hơn và có thể kiểm tra lại sau khi đợt phân bổ kết thúc.

% =====================================================================
\section{Giải pháp kiến trúc phân bổ linh hoạt theo năm học}
\label{sec:contrib_policy}
% =====================================================================

\subsection{Dẫn dắt và bài toán}
\label{subsec:policy_problem}

\subsubsection{Vấn đề cần giải quyết}

Một yêu cầu phát sinh trong quá trình demo là: ``Năm sau muốn tăng quota năm
nhất từ 40\% lên 45\% thì làm thế nào?'' Nếu các tham số quota vẫn được
hard-code trong mã nguồn, mỗi thay đổi như vậy sẽ cần lập trình viên sửa code,
kiểm thử và triển khai lại. Yêu cầu này cho thấy phần cấu hình chính sách
không nên được xem như hằng số kỹ thuật, mà cần trở thành dữ liệu nghiệp vụ có
thể quản lý qua hệ thống.

Chính sách phân bổ ký túc xá không cố định mà thay đổi theo từng năm học, thậm
chí có thể cần điều chỉnh giữa kỳ khi chỉ tiêu, cơ cấu sinh viên hoặc quy định
vận hành thay đổi. Tỷ lệ ưu tiên theo nhóm năm có thể điều chỉnh, tiêu chí xét
duyệt có thể mở rộng, ngưỡng khoảng cách có thể thay đổi. Trong hệ thống với
logic viết cứng trong mã nguồn, mỗi thay đổi chính sách đòi hỏi một chu kỳ
phát triển phần mềm mới --- độ trễ này không phù hợp với nhịp thay đổi của
công tác quản lý ký túc xá.

Bên cạnh đó, cơ cấu sinh viên theo nhóm năm dịch chuyển mỗi năm (cohort
shift): khi K69 lên năm 2, K70 vào thay vị trí năm 1. Nếu không có công cụ
dự báo, cán bộ buộc phải phân bổ chỉ tiêu dựa trên kinh nghiệm chủ quan thay
vì dữ liệu thực tế. Đây là bài toán lập kế hoạch chủ động, không phải phản
ứng sau khi vấn đề đã xảy ra.

\subsubsection{Đối tượng hưởng lợi}

\textit{Cán bộ quản lý:} có thể tự cập nhật quota và quy định phân bổ theo
năm học hoặc giữa kỳ qua giao diện cấu hình mà không cần yêu cầu sửa mã nguồn.
\textit{Ban lãnh đạo ký túc xá:} có lịch sử chính sách đầy đủ theo từng phiên
bản và từng năm học để đối chiếu, kiểm toán. \textit{Toàn bộ sinh viên:} được
hưởng lợi từ chính sách phân bổ phản ánh sát hơn thực tế cơ cấu năm học hiện
tại.

\subsubsection{Điểm khác biệt so với giải pháp hiện có}

\begin{table}[H]
\centering
\caption{So sánh giải pháp quản lý chính sách phân bổ}
\label{tab:policy_comparison}
\small
\begin{tabular}{|p{3.5cm}|p{4.5cm}|p{4.8cm}|}
\hline
\textbf{Tiêu chí} & \textbf{Logic hard-coded hiện có} & \textbf{Policy-Driven Framework} \\
\hline
Thay đổi chính sách & Cần chu kỳ phát triển phần mềm & Cấu hình qua giao diện, vài phút \\
\hline
Kế thừa năm trước & Không có, cấu hình lại từ đầu & Tự động kế thừa, admin chỉ điều chỉnh \\
\hline
Lịch sử chính sách & Không được lưu trữ & Lưu đầy đủ theo phiên bản và năm học \\
\hline
Dự báo cohort shift & Không có & Module tự động dựa trên dữ liệu lịch sử \\
\hline
Tách biệt quota và điểm & Không (lẫn lộn trong code) & Tách thành hai thực thể độc lập \\
\hline
\end{tabular}
\end{table}

\subsection{Giải pháp}
\label{subsec:policy_arch}

Lúc đầu, \texttt{QuotaConfig} được lồng trực tiếp bên trong
\texttt{AllocationPolicy} vì hai khái niệm này cùng liên quan đến phân bổ
phòng. Tuy nhiên, yêu cầu điều chỉnh tỷ lệ quota mà không thay đổi công thức
tính điểm ưu tiên cho thấy hai thành phần này có vòng đời khác nhau: quota
thay đổi theo cơ cấu sinh viên từng năm hoặc từng đợt, còn trọng số điểm ưu
tiên có thể giữ nguyên nhiều năm liền. Từ đó, đồ án hình thành kiến trúc bốn
lớp để tách cấu hình quota, chính sách tính điểm, chu kỳ phân bổ và dịch vụ
thực thi.

Hình~\ref{fig:policy_driven_allocation} trình bày kiến trúc Policy-Driven
Allocation Framework: chính sách phân bổ được lưu trữ dưới dạng tài liệu có
cấu trúc trong cơ sở dữ liệu, hoàn toàn tách khỏi mã nguồn.

\begin{figure}[H]
\centering
\includegraphics[width=0.85\textwidth]{Hinhve/policy_driven_allocation.png}
\caption{Kiến trúc quản lý chính sách phân bổ động}
\label{fig:policy_driven_allocation}
\end{figure}

Kiến trúc gồm bốn lớp tương tác theo chiều dọc. \texttt{QuotaConfig} định
nghĩa tỷ lệ chỉ tiêu theo nhóm năm --- thực thể thay đổi hàng năm theo biến
động cơ cấu sinh viên. \texttt{AllocationPolicy} định nghĩa trọng số tính
điểm ưu tiên và tham chiếu về \texttt{QuotaConfig} --- thực thể ổn định hơn,
có thể dùng qua nhiều năm. \texttt{AllocationCycle} xác định khoảng thời gian
đăng ký và liên kết với một chính sách cụ thể. \texttt{AllocationService} đọc
cấu hình từ ba lớp trên và thực thi phân bổ mà không chứa bất kỳ giá trị
hard-coded nào.

Ở tầng dữ liệu, \texttt{QuotaConfig} được quản lý theo trạng thái bản nháp và
bản đã công bố, kèm các trường \texttt{version}, \texttt{isDraft},
\texttt{isPublished}, \texttt{publishedAt} và \texttt{publishedBy}. Khi một
bản quota được công bố, dịch vụ \texttt{quotaPublishService} cập nhật quota
đang có hiệu lực, tăng phiên bản và tạo hoặc cập nhật \texttt{AllocationPolicy}
tương ứng thông qua \texttt{sourceQuotaId}. Cách tổ chức này giúp hệ thống lưu
được lịch sử thay đổi để đối chiếu về sau.

Quyết định thiết kế then chốt là tách \texttt{QuotaConfig} khỏi
\texttt{AllocationPolicy}: hai khái niệm nghiệp vụ có vòng đời thay đổi khác
nhau không nên lưu trong cùng một cấu trúc dữ liệu. Nhờ tách biệt này, cán bộ
có thể điều chỉnh tỷ lệ chỉ tiêu mà không ảnh hưởng đến công thức tính điểm
ưu tiên, và ngược lại.

Công thức theo dõi mức độ tuân thủ quota sau mỗi đợt phân bổ:
\begin{equation}
\Delta_g = n_g^{\mathrm{thực\ tế}} - \lfloor Q_g \cdot C_{\mathrm{tổng}} \rfloor
\label{eq:quota_deviation}
\end{equation}
Khi $\Delta_g > 0$, nhóm $g$ đang vượt chỉ tiêu; $\Delta_g < 0$ là chưa đạt.
Hệ thống hiển thị chỉ số này theo thời gian thực để cán bộ điều chỉnh cho các
đợt đăng ký tiếp theo.


\subsection{Kết quả đạt được}
\label{subsec:policy_result}

Kiến trúc Policy-Driven Allocation mang lại ba giá trị định lượng được.
\textit{Thứ nhất}, thời gian cập nhật chính sách mỗi năm hoặc giữa kỳ giảm từ
vài ngày (phát triển phần mềm mới) xuống còn vài phút (điền form cấu hình).
\textit{Thứ hai}, lịch sử chính sách được bảo toàn theo năm học và phiên bản
thông qua trường \texttt{version} và \texttt{sourceQuotaId}, hỗ trợ kiểm toán
và đối chiếu.
\textit{Thứ ba}, module dự báo cohort shift đề xuất quota tự động dựa trên
xu hướng $k$ năm gần nhất:
\begin{equation}
\hat{Q}_g^{\mathrm{đề\ xuất}} = \frac{1}{k} \sum_{i=1}^{k}
    \frac{n_{g,i}^{\mathrm{thực\ tế}}}{\displaystyle\sum_{g'} n_{g',i}^{\mathrm{thực\ tế}}}
\label{eq:quota_recommendation}
\end{equation}
cho phép cán bộ lập kế hoạch chủ động trước khi năm học mới bắt đầu thay vì
phản ứng sau khi vấn đề đã xảy ra.

\subsection{Giá trị mang lại}
\label{subsec:policy_value}

Giá trị của kiến trúc chính sách nằm ở việc chuyển quota và các quy định phân
bổ từ mã nguồn sang dữ liệu cấu hình có vòng đời riêng. Cán bộ có thể thay đổi
tỷ lệ chỉ tiêu, thời gian hiệu lực, quy định vượt chỉ tiêu hoặc cửa sổ dữ liệu
phân tích theo từng năm học, thậm chí giữa kỳ, mà không cần sửa mã nguồn và
không phụ thuộc vào một chu kỳ triển khai phần mềm mới.

Việc lưu phiên bản và lịch sử công bố giúp mỗi thay đổi có dấu vết rõ ràng:
quota nào đang có hiệu lực, ai công bố, công bố vào thời điểm nào và chính
sách phân bổ nào được sinh ra từ bản quota đó. Khi cần đối chiếu kết quả phân
bổ của một năm học hoặc giải trình về thay đổi chỉ tiêu, hệ thống có thể truy
ngược lại cấu hình đã được sử dụng thay vì chỉ dựa vào ghi nhớ hoặc tài liệu
rời.

Tách cấu hình khỏi mã nguồn cũng làm ranh giới trách nhiệm rõ hơn. Cán bộ quản
lý chịu trách nhiệm điều chỉnh quy định vận hành trong phạm vi được giao; phần
mềm chịu trách nhiệm kiểm tra tính hợp lệ, lưu phiên bản và thực thi nhất quán
theo cấu hình đã công bố.

% =====================================================================
\section{Tổng kết chương}
\label{sec:contrib_summary}
% =====================================================================

Ba đóng góp được trình bày trong chương này tuy giải quyết các bài toán khác
nhau nhưng có cùng định hướng: biến các nhu cầu vận hành cụ thể của ký túc xá
thành những cơ chế kỹ thuật có thể triển khai, kiểm tra và mở rộng. Trực quan
hóa phòng giúp sinh viên có căn cứ quan sát rõ hơn trước khi đăng ký; phân bổ
phòng dựa trên quota giúp giảm thao tác thủ công và tạo kết quả có thể truy
vết; kiến trúc chính sách giúp cán bộ thay đổi quy định theo năm học hoặc giữa
kỳ mà không cần sửa mã nguồn.

Đối với \textbf{đóng góp thứ nhất} (mục~\ref{sec:contrib_visualization}), đồ án
đã phát triển giải pháp từ ảnh tĩnh sang hai chế độ trực quan hóa gồm Quick
Tour 360° và Free Explore 3D. Điểm đáng chú ý không chỉ là việc dựng được mô
hình 3D, mà còn là quá trình tối ưu để đưa kích thước mô hình từ hơn 200~MB
xuống dưới 8~MB và thiết kế cơ chế tải phù hợp với năng lực thiết bị. Nhờ đó,
tính năng có khả năng sử dụng thực tế trên môi trường web và hỗ trợ tốt hơn cho
việc giới thiệu phòng, ký túc xá.

Đối với \textbf{đóng góp thứ hai} (mục~\ref{sec:contrib_allocation}), đồ án đã
tách quy trình phân bổ thành các bước rõ ràng gồm kiểm tra ràng buộc bắt buộc,
xếp hạng theo điểm ưu tiên, gán phòng và ghi nhật ký. Cách tổ chức này giúp
quy trình phân bổ khách quan hơn, nhất quán hơn và có căn cứ đối chiếu, đồng
thời giảm phụ thuộc vào Excel và các thao tác thủ công dễ phát sinh sai sót.

Đối với \textbf{đóng góp thứ ba} (mục~\ref{sec:contrib_policy}), đồ án đã tách
quota, chính sách tính điểm, chu kỳ phân bổ và dịch vụ thực thi thành các lớp
có trách nhiệm riêng. Nhờ cơ chế bản nháp, công bố, phiên bản và tham chiếu
\texttt{sourceQuotaId}, các thay đổi chính sách được lưu lại để kiểm tra về
sau, còn cán bộ có thể điều chỉnh quota hoặc quy định vận hành mà không cần
can thiệp vào mã nguồn.

Như vậy, Chương 5 đã làm rõ phần đóng góp kỹ thuật của đồ án trên ba khía
cạnh: trải nghiệm ra quyết định của sinh viên, quy trình phân bổ phòng và khả
năng thích ứng của chính sách vận hành. Từ các đóng góp này, Chương 6 sẽ tiếp
tục đánh giá kết quả đạt được của toàn hệ thống, chỉ ra những hạn chế còn tồn
tại và đề xuất các hướng phát triển tiếp theo.


\end{document}
